\subsection{\texorpdfstring{Hypocontinuity of the mapping \((u, v) \mapsto v \circ u\)}{Hypocontinuity of the mapping (u, v) \textbackslash mapsto v \textbackslash circ u}}

PROPOSITION 9. --- Let R, S, T be three locally convex Hausdorff spaces. Suppose that the spaces \(\mathcal{L}(\mathbf{R};\mathbf{S})\) , \(\mathcal{L}(\mathbf{S};\mathbf{T})\) , \(\mathcal{L}(\mathbf{R};\mathbf{T})\) are each assigned the topology of simple (resp. compact, bounded) convergence. Then the bilinear mapping \((u, v) \mapsto v \circ u\) from \(\mathcal{L}(\mathrm{R}; \mathrm{S}) \times \mathcal{L}(\mathrm{S}; \mathrm{T})\) into \(\mathcal{L}(\mathrm{R}; \mathrm{T})\) is \((\mathfrak{S}, \mathfrak{T})\) -hypocontinuous, where \(\mathfrak{T}\) is the family of equicontinuous subsets of \(\mathcal{L}(\mathrm{S}; \mathrm{T})\) , and \(\mathfrak{S}\) the family of finite (resp. compact, bounded) subsets of \(\mathcal{L}(\mathrm{R}; \mathrm{S})\) .

We first prove that \((u, v) \mapsto v \circ u\) is \(\mathfrak{T}\) -hypocontinuous. Let H be an equicontinuous set in \(\mathcal{L}(S; T)\) , let W be a neighbourhood of 0 in T and let M be a finite (resp. compact, bounded) subset of R. We must show that there exists a neighbourhood V of 0 in S such that if \(u(M) \subset V\) and \(v \in H\) , then \(v(u(M)) \subset W\) . But for this, it is enough to have \(v(V) \subset W\) for all \(v \in H\) , and the existence of such a neighbourhood follows from the equicontinuity of H.

To see that \((u, v) \mapsto v \circ u\) is \(\mathfrak{S}\) -hypocontinuous, we shall prove that, for every neighbourhood \(W\) of 0 in \(T\) , every finite (resp. compact, bounded) subset \(M\) of \(R\) and every finite (resp. compact, bounded) subset \(L\) of \(\mathcal{L}(R; S)\) there exists a finite (resp. compact, bounded) subset \(N\) of \(S\) such that the relations \(v(N) \subset W\) and \(u \in L\) imply that \(v(u(M)) \subset W\) . Evidently it is enough to show that we can take \(N = \bigcup_{u \in L} u(M)\) ,

i.e.~that the set N is finite (resp. compact, bounded) whenever L and M are. This is immediate if L and M are finite, or if M is bounded in R and L is bounded in \(\mathcal{L}(\mathrm{R};\mathrm{S})\) (for the topology of bounded convergence, cf.~III, p.~22). Finally, we show that if M is compact in R and L is compact in \(\mathcal{L}(\mathrm{R};\mathrm{S})\) for the topology of compact convergence, then N is compact in S. But if \(u_{M}\) is the restriction to M of \(u \in L\) , the mapping \(u \mapsto u_{M}\) from L into the space \(\mathcal{C}(\mathrm{M};\mathrm{S})\) of all continuous mappings from M into S, with the topology of uniform convergence, is continuous; hence the image of L under this mapping is compact, and our assertion then follows from the continuity of the map \((w,x) \mapsto w(x)\) from \(\mathcal{C}(\mathrm{M};\mathrm{S}) \times \mathrm{M}\) into S (GT, X, § 1, No.~6, prop. 9).

In the two corollaries that follow, we assume as in prop. 9, that the spaces \(\mathcal{L}(\mathbf{R};\mathbf{S})\) , \(\mathcal{L}(\mathbf{S};\mathbf{T})\) , \(\mathcal{L}(\mathbf{R};\mathbf{T})\) are all three assigned the topology of simple convergence, or all three the topology of compact convergence, or all three that of bounded convergence.

COROLLARY 1. --- For every equicontinuous subset H of \(\mathcal{L}(\mathbf{S};\mathbf{T})\) the map \((u,v)\mapsto v\circ u\) from \(\mathcal{L}(\mathbf{R};\mathbf{S})\times \mathbf{H}\) into \(\mathcal{L}(\mathbf{R};\mathbf{T})\) is continuous.

This follows immediately from prop. 9 (III, p.~32) and 4 (III, p.~31).

COROLLARY 2. --- Suppose S is barrelled. If the sequence \((u_n)\) tends to \(u\) in \(\mathcal{L}(\mathbf{R};\mathbf{S})\) and the sequence \((v_{n})\) to \(v\) in \(\mathcal{L}(\mathbf{S},\mathbf{T})\) , then the sequence \((v_{n}\circ u_{n})\) tends to \(v\circ u\) in \(\mathcal{L}(\mathbf{R};\mathbf{T})\) .

In fact, the sequence \((v_{n})\) , being simply bounded in \(\mathcal{L}(\mathrm{S};\mathrm{T})\) is equicontinuous, since \(\mathrm{S}\) is barrelled (III, p.~25, th. 1); the corollary is then a consequence of cor. 1.

